\section{Implementation}
\label{sec:implementation}

\subsection{Components Overview}

Our prototype extends Linux 6.15 with a small core \sys kernel code source file ($\sim$1 KLOC) and a lightweight integration into NVIDIA's open GPU kernel modules ($\sim$100 LOC changed across UVM and scheduler code). These changes register new eBPF program types, hook contexts, and kfuncs, and wire struct\_ops tables into the GPU memory and scheduling paths (§4.1–§4.2). The userspace control plane and loader/JIT are substantially larger ($\sim$10 KLOC), implemented by extending bpftime's loader infrastructure. Finally, our eBPF$\rightarrow$PTX/SPIR‑V backend, built on LLVM, consists of $\sim$1 KLOC of lowering code that translates the restricted \sys eBPF subset into device code.

\subsection{Host-side Integration}
\label{sec:impl-host}

\subsubsection{Driver hooks for memory and scheduling}

\sys builds on the NVIDIA Open GPU Kernel Modules (Turing, Ampere, Hopper and
newer) to avoid binary patching and to track vendor updates. In the UVM module,
we insert callbacks at critical memory decision points: page fault handlers,
prefetch logic, and reclamation paths. These implement the region-level
Add/Access/Remove and prefetch events from Section~\ref{sec:design-mem}. Each
callback constructs a compact context (for example region identifier or fault
address, region size, process and tenant IDs, and local memory pressure) and
invokes an attached eBPF handler. The handler consults eBPF maps that hold
policy state (LRU or LFU metadata, tenant budgets, classification flags) and
returns compact decisions such as admit, bypass, pin, evict, or a ranking over
candidate regions. The driver applies these choices using its existing mechanisms
for migration and mapping updates.

For scheduling, we extend the kernel lifecycle code that manages GPU contexts and
hardware queues with struct\_ops callback vectors at kernel submission, admission,
preemption, and completion, implementing the interface in
Section~\ref{sec:design-sched}. Submission callbacks see kernel configuration
(grid and block size), tenant or stream IDs, and priority hints. Admission
callbacks run before dispatch and can re-order ready queues or delay admission.
Preemption callbacks fire at cooperative yield points or at scheduling quanta and
allow policies to decide whether to continue or preempt a kernel based on its
progress and competing work. Completion callbacks update per-tenant statistics
and virtual time. All callbacks can be attached and detached dynamically without
driver recompilation or node restart.

On the host, \sys reuses the standard kernel eBPF verifier to ensure memory
safety, bounded loops, stack limits, and controlled helper and map usage. We implement
a kernel module that extends the verification context to support our GPU-specific attach points
and special kfuncs. This module leverages the kernel's annotation mechanisms (BTF and kfunc
declarations) to enforce interface correctness and type safety, similar to the approach
in the bpftime paper~\cite{bpftime}. Where policies must trigger specific actions in the driver
(such as committing a victim selection or marking a kernel for preemption), we expose these
interfaces as safe kfuncs that can be called from eBPF under verifier control.

\subsubsection{\sysdevice user-space interface}

\sysdevice provides eBPF-style uprobe and uretprobe hooks for GPU-related
user-level functions, enabling bcc-style observability tools and higher-level
policy logic without modifying applications. Developers attach eBPF programs to
function entry and exit to collect metrics such as per-call latency, counts, or
arguments, using standard eBPF maps.

At kernel launch time, \sys optionally JIT-translates these handlers to device
code and injects entry and exit trampolines so that the same logic can run on
the GPU, recording per-warp timestamps or counters directly into shared maps.
This lifts the Linux uprobe semantics onto devices while retaining the ability
to enable or disable instrumentation at runtime without relaunching applications.

\subsection{GPU Runtime}
\label{sec:impl-gpu}

\subsubsection{Code generation and trampolines}
gBPF builds its GPU compiler on bpftime’s loader front end: when a gBPF program is verified and loaded into the kernel, the bpftime-based loader intercepts it, extracts the bytecode and BTF metadata, and hands them to a user-space backend (LLVM-based) that lowers the restricted eBPF subset to device code (PTX today, SPIR-V or other ISAs in principle). We disable bpftime’s user-space eBPF runtime; all execution remains in the kernel (via the normal eBPF JIT) or on the GPU via our backend and runtime. To avoid recompiling applications, gBPF patches kernels at launch with small trampolines at entry, sampled memory operations, and phase boundaries: each trampoline saves minimal live state into per-warp scratch, builds a compact context, runs the translated handler in a single designated lane, and restores state. This design treats verifier-accepted eBPF as the single policy IR and carries it unchanged from driver hooks into GPU kernels via trampolines, rather than introducing a separate GPU-only DSL or ad hoc SASS rewriter, the only difference is the capabilities of the API and restrictions of the gBPF subset.



\subsubsection{Sandbox and resource bounding}
All device-side gBPF code executes inside a lightweight sandbox that mirrors the verifier’s assumptions and adds GPU-specific execution discipline. At load time, gBPF queries verifier metadata (instruction count, stack depth, register usage) and uses it to set conservative per-invocation work budgets and sampling rates for each hook. The backend restricts all data-dependent memory operations to gBPF-managed maps and per-warp scratch, emitting explicit bounds checks against known sizes, and never generates GPU barriers, queue operations, or raw pointers into application or driver memory; only a small set of runtime helpers with verifier-approved prototypes is exposed. Each translated handler is wrapped in a minimal prologue/epilogue that saves and restores the designated lane’s registers and predicates and runs the policy logic under these guards. Compared to existing GPU instrumentation frameworks that execute arbitrary PTX/SASS without static guarantees, gBPF is the first to tie device-side policy execution to the Linux eBPF verifier and to a SIMT-aware sandbox, providing verifier-backed memory safety and bounded, policy-controlled work for GPU-resident extensions.


\subsubsection{SIMT-specific optimizations}

To execute scalar policy logic efficiently on a SIMT architecture, \sys aggregates
execution at warp granularity. When a device hook fires, a designated lane (typically
lane 0) in the warp executes the eBPF handler, while other lanes either remain idle
or contribute minimal input via warp-level primitives such as shuffle instructions.
This avoids intra-warp divergence inside policy code and reduces redundant work.

Updates to global eBPF maps use warp-local and SM-local accumulators in registers
or shared memory, which are periodically flushed to global memory, reducing atomic
contention and global traffic. Context structures passed to handlers are minimized
based on static analysis: only fields that may be read are materialized, and
frequently used identifiers are precomputed in the trampoline. Configurable sampling
counters implemented in registers let handlers run only every \(N\)-th event; sampling
rates are stored in control-plane maps and can be tuned at runtime.

For scheduling policies, we implement a thread-block scheduler that models each
kernel as a collection of work units (for example tiles or logical block indices)
consumed by persistent worker blocks. On recent NVIDIA GPUs we map this to cluster
launch control, which treats a grid of blocks as cancellable work units. Each worker
processes its current unit, invokes a device-side \sys handler (running in one lane)
to select or approve the next work unit based on shared maps encoding per-tenant
priorities, global queue state, and per-block progress, and uses the cluster API to
claim or release indices. Progress and behavior probes implemented with \sys
trampolines at kernel entry and annotated phase boundaries record per-warp or
per-block iteration counts, phase identifiers, and simple classification signals
(such as memory-bound versus compute-bound phases) into shared maps that are
visible to both device and host schedulers. All device-side scheduling logic is
subject to the same sandbox and bounded-execution constraints as other \sys code.

\subsection{Cross-domain State and Synchronization}
\label{sec:impl-state}

\subsubsection{Shared maps and allocation}

\sys uses eBPF maps as the unified state abstraction across host and device. We
allocate maps and coordination buffers in memory regions accessible to CPU and GPU,
avoiding explicit marshalling. On our NVIDIA prototypes we use managed shared
memory so that map contents are addressable in both contexts. Maps store policy
state such as per-region hotness, per-tenant scheduling priorities, per-block
progress, and counters for observability tools.

Within each domain, we use standard atomics and interprocess primitives on the
host and CUDA atomics and fences on the device over well-defined layouts in
shared memory. To reduce contention, only a subset of threads (for example one
lane per warp) performs atomic updates while other threads aggregate locally in
registers or shared memory, and hot maps can be sharded by key when necessary.

\subsubsection{CPU–GPU visibility and fences}

Because CPU and GPU execute independently, cross-domain consistency requires
explicit fences. To ensure GPU writes become visible to the host, device code
calls a small wrapper around \texttt{\_\_threadfence\_system}:

\begin{verbatim}
__device__ inline void memsys_bar_gpu() {
  __threadfence_system();
}
\end{verbatim}

To ensure host writes are visible to the GPU, host code calls a companion wrapper
around \texttt{cudaStreamSynchronize}:

\begin{verbatim}
static inline void memsys_bar_cpu(cudaStream_t stream = 0) {
  cudaStreamSynchronize(stream);
}
\end{verbatim}

Persistent kernels invoke \texttt{memsys\_bar\_gpu()} before reading state that the
host may have updated; host code invokes \texttt{memsys\_bar\_cpu()} after writing
maps that device-side handlers will consume. These abstractions centralize tuning
and simplify reasoning about CPU–GPU synchronization for \sys policies.

